% Treasury bill cash-flow timeline (L2a).
% Keep this figure structural: the valuation and return derivations belong in
% the surrounding notebook and slides.
% Flow convention: an arrow leaving a time-point unit is money paid;
% an arrow entering one is money received.
\documentclass[tikz,border=6pt]{standalone}
\usepackage{vnfigure}

\begin{document}
\begin{tikzpicture}[x=2.0cm,y=1cm,>=Latex,font=\sffamily\small]

  % Span label ------------------------------------------------
  \draw[vn brace] (0.45,1.65) -- (3.40,1.65)
    node[midway,above=7pt,text=vnink]
    {\bfseries zero-coupon bill: one payment, one receipt};

  % Timeline --------------------------------------------------
  \draw[vn timeline] (0,0) -- (4.05,0) node[right] {time};
  \node[vn unit] (b0) at (0.45,0) {};
  \node[vn unit] (bT) at (3.40,0) {};
  \node[below=4pt] at (b0) {$t=0$};
  \node[below=4pt] at (bT) {$t=T$};

  % Purchase price paid at the valuation date ------------------
  \draw[vn paid] (b0.north) -- (0.45,1.05);
  \node[right,text=vncarnelian] at (0.45,0.62) {$-V_{B}$};
  \node[above,vn annotation] at (0.45,1.05) {cash paid};

  % Par value received at maturity -----------------------------
  \draw[vn received] (3.40,1.05) -- (bT.north);
  \node[right,text=vnlink] at (3.40,0.62) {$+V_{P}$};
  \node[above,vn annotation] at (3.40,1.05) {cash received};

\end{tikzpicture}
\end{document}
